\begin{definition}[The program of one process]
  \label{def:aba-programs}
  \lean{PLTS.ABA.GBCA.ByABDY.RoundRecord}\lean{PLTS.ABA.RoundLoopRecord}\lean{PLTS.ABA.ABDY.RoundRecordMap}\lean{PLTS.ABA.ABDY.ProcessRecord}\lean{PLTS.ABA.Composition.NetworkEvent}\lean{PLTS.ABA.Composition.ExtendedLabel}\lean{PLTS.ABA.Implementation.networkEventLabels}\lean{PLTS.ABA.ABDY.RoundStep}\lean{PLTS.ABA.ABDY.ABAProgramStep}\lean{PLTS.ABA.ABDY.ABAProgram}\leanok
  \uses{def:aba-labels, def:aba-gbca-vocab, def:aba-core-vocab, def:aba-implementation, def:relabel, def:syncproduct}
  The program of process $j$ of the protocol, Definition~\ref{def:aba-implementation} at ABDY22's implementation of Definition~\ref{def:aba-gbca-vocab}, on the state $\mathrm{ProcessRecord}(n) = \mathrm{RoundLoopRecord}(n) \times \mathrm{RoundRecordMap}(n)$ --- the round-loop record beside the round record of every round the process has touched and its termination flag, which the round advance does not reset and which the row $terminate$ leaves standing (D22): \leandecl{PLTS.ABA.ABDY.ABAProgram}.
  The round's own rows are the fourteen of \leandecl{PLTS.ABA.ABDY.RoundStep}: the eight multicasts of those levels, the delivery, the two calls and the three returns.
  It reads no copy of the corrupted set and no budget, only the replacement flag \leandecl{PLTS.ABA.RoundLoopRecord.corrupted} of its own record, which \leandeclnamed{failSelf}{PLTS.ABA.Implementation.ProgramStep.failSelf} writes and every other row requires to be down; in place of those rows the replaced program is the self-loop \leandeclnamed{corruptedIdle}{PLTS.ABA.Implementation.ProgramStep.corruptedIdle} (D23).
\end{definition}
